\chapter{Myasthenia Gravis}

\medskip
\noindent\textit{\textbf{Under review.} This chapter is an AI-assisted educational draft; it carries no source citations and is not a substitute for professional medical advice.}

\section*{Definition and Overview}
Myasthenia gravis (MG) is a chronic autoimmune neuromuscular disorder characterised by fluctuating weakness of skeletal (voluntary) muscles, typically worse with activity and improving with rest. It is caused by autoantibodies that disrupt transmission at the neuromuscular junction (NMJ), most commonly antibodies against the acetylcholine receptor (AChR) but also against muscle-specific kinase (MuSK), LRP4, or agrin. The weakness most often affects the eye muscles (ocular MG) and may generalise (generalised MG) to involve bulbar, limb, and respiratory muscles. The disease can range from mild, isolated eye weakness to life-threatening respiratory failure (myasthenic crisis). Onset and severity vary; treatment can achieve good control for most patients.

\section*{Epidemiology}
Myasthenia gravis affects approximately 15--20 per 100,000 people; in Australia this equates to roughly 4,500--5,000 people living with the condition. It has a bimodal age distribution: early-onset MG (more common in women, typically 20--40 years) and late-onset MG (more common in men, typically after 50 years). Ocular MG is more common in young men and in later-onset patients. Roughly 10--15\% of MG is associated with a thymoma, and thymic hyperplasia is present in many younger patients. Prevalence has been gradually increasing due to improved diagnosis and survival, as well as ageing of the population.

\section*{Aetiology and Pathophysiology}
\begin{itemize}
    \item \textbf{Autoimmunity:} production of autoantibodies targeting proteins of the neuromuscular junction (NMJ)
    \item \textbf{AChR antibodies:} present in ~80--85\% of generalised MG and ~50\% of ocular MG; block or destroy AChRs, reducing postsynaptic response (complement-mediated damage)
    \item \textbf{MuSK antibodies:} present in a minority (often AChR-negative) and associated with prominent bulbar and respiratory weakness
    \item \textbf{Thymus:} abnormal thymus (hyperplasia or thymoma) plays a central role; thymectomy can improve outcome, especially in AChR-positive, early-onset disease
    \item \textbf{Pathogenesis:} autoantibodies reduce the number of functioning AChRs at the NMJ, so repetitive nerve stimulation depletes available receptors and causes fatigable weakness
    \item \textbf{Triggers:} infections, stress, some medications (e.g., aminoglycosides, beta-blockers, certain anaesthetics), surgery, and pregnancy can worsen symptoms
\end{itemize}

\section*{Clinical Features}
\begin{itemize}
    \item \textbf{Ocular symptoms:} ptosis (drooping eyelid) and diplopia (double vision), the presenting feature in the majority and often fluctuating
    \item \textbf{Fatigable weakness:} weakness increases with repeated use and improves after rest (fatiguability)
    \item \textbf{Bulbar symptoms:} difficulty chewing, swallowing (dysphagia), slurred speech (dysarthria), and facial weakness
    \item \textbf{Limb weakness:} proximal muscles (shoulders, hips) more than distal, affecting arms, legs, and neck (head drop)
    \item \textbf{Respiratory weakness:} shortness of breath on exertion; severe involvement causes respiratory failure (myasthenic crisis)
    \item \textbf{Diurnal variation:} symptoms often worse later in the day or with activity
    \item \textbf{Absent sensory features:} sensation, reflexes, and coordination are usually normal
\end{itemize}

\section*{Differential Diagnosis}
\begin{itemize}
    \item Ocular: caused by other conditions such as third nerve palsy, internuclear ophthalmoplegia, ophthalmoplegic migraine, or thyroid eye disease
    \item Lambert-Eaton myasthenic syndrome (LEMS): presynaptic disorder, often paraneoplastic, with different electrophysiology and autonomic symptoms
    \item Congenital myasthenic syndromes
    \item Botulism: acute descending weakness with autonomic features
    \item Motor neurone disease, myopathies (e.g., inclusion body myositis), or periodic paralysis
    \item Stroke, brainstem lesion, or psychogenic weakness
    \item Guillain-Barré syndrome (if acute and generalised)
\end{itemize}

\section*{When to Seek Medical Care}
Seek medical assessment for new or worsening muscle weakness, especially fatigable symptoms, double vision, drooping eyelids, or difficulty swallowing or breathing. Early referral to a neurologist is important for diagnosis and treatment.

\textbf{Call 000 or attend an emergency department for (myasthenic crisis):}
\begin{itemize}
    \item Severe difficulty breathing or shortness of breath
    \item Inability to speak, swallow, or clear secretions
    \item Sudden severe generalised weakness
    \item Signs of respiratory distress or impending respiratory failure
\end{itemize}

\section*{Diagnosis and Investigations}
\begin{itemize}
    \item \textbf{Antibody testing:} serum AChR antibodies; if negative, MuSK, LRP4, and agrin antibodies
    \item \textbf{Electrophysiology:} repetitive nerve stimulation showing decremental response; single-fibre EMG (very sensitive)
    \item \textbf{Edrophonium (Tensilon) test:} less commonly used now; gives transient improvement but carries cardiac risk
    \item \textbf{Chest imaging:} CT or MRI of the chest/mediastinum to detect thymoma
    \item \textbf{Thyroid function:} exclude thyroid disease (thyrotoxicosis can mimic/worsen weakness)
    \item \textbf{Other:} exclude coexisting autoimmune conditions (e.g., vitiligo, rheumatoid arthritis, lupus)
\end{itemize}

\section*{Management}
\begin{itemize}
    \item \textbf{Acetylcholinesterase inhibitors:} pyridostigmine (e.g., 30--60 mg every 4--6 hours) for symptomatic treatment
    \item \textbf{Immunosuppression:} corticosteroids (prednisolone) and steroid-sparing agents (azathioprine, mycophenolate, tacrolimus)
    \item \textbf{Thymectomy:} surgical removal of the thymus, particularly for AChR-positive, early-onset generalised MG; essential if thymoma
    \item \textbf{Plasma exchange / IVIG:} for acute exacerbations or pre-operatively; rapidly reduce circulating antibodies
    \item \textbf{Novel therapies:} complement inhibitors (eculizumab) and FcRn antagonists for refractory AChR-positive MG
    \item \textbf{Avoidance of triggers:} avoid medications that worsen MG (e.g., aminoglycosides, quinolones, magnesium, some beta-blockers); treat infections promptly
    \item \textbf{Multidisciplinary care:} neurologist, respiratory support, speech therapy for swallowing, and ophthalmology input
    \item \textbf{Myasthenic crisis:} intensive care with respiratory support, plasma exchange or IVIG, and treatment of the precipitating cause
\end{itemize}

\section*{Complications}
\begin{itemize}
    \item Myasthenic crisis: severe respiratory weakness requiring intubation and ventilation (a medical emergency)
    \item Cholinergic crisis from excess pyridostigmine (muscarinic side effects: diarrhoea, salivation, bradycardia)
    \item Treatment-related complications: corticosteroid side effects, immunosuppression-related infections
    \item Dysphagia leading to aspiration pneumonia and malnutrition
    \item Thymoma (present in up to 15\%): requires resection and ongoing surveillance
    \item Effects on pregnancy (fluctuating), with need for close specialist monitoring
\end{itemize}

\section*{Prognosis and Outcomes}
With modern treatment, the prognosis for myasthenia gravis is excellent for most patients: the majority achieve good to excellent control of symptoms. Disease activity tends to fluctuate, with the worst weakness typically in the first few years, then often stabilising or improving. Treatment-related remission or marked improvement occurs in many patients on immunosuppression; some enter long remissions. Ocular MG may remain localised or progress to generalised MG (more likely within the first 2--3 years). Myasthenic crises are now associated with much lower mortality (~5\%) with modern intensive care. Life expectancy approaches normal with treatment.

\section*{Prevention and Self-Care}
\begin{itemize}
    \item \textbf{Medication awareness:} inform all health professionals of the diagnosis; check drug interactions before taking new medicines
    \item \textbf{Infection prevention:} prompt treatment of infections and up-to-date vaccinations (with specialist advice)
    \item \textbf{Energy conservation:} plan activities to manage fatiguability and rest when needed
    \item \textbf{Safety:} avoid driving or operating machinery during episodes of significant weakness or diplopia, with medical advice
    \item \textbf{Adherence:} take prescribed immunosuppressants and anticholinesterases as directed; do not stop suddenly
    \item \textbf{Emergency plan:} know the warning signs of crisis and how to seek urgent care
    \item \textbf{Support:} patient organisations (MG Association of Australia) provide education and peer support
\end{itemize}

\section*{Sources and Further Reading}
\begin{itemize}
    \item Myasthenia Gravis Association of Australia. https://myasthenia-gravis.org.au/
    \item Healthdirect Australia. \textit{Myasthenia gravis.} https://www.healthdirect.gov.au/
    \item National Institute of Neurological Disorders and Stroke (NINDS). \textit{Myasthenia Gravis Fact Sheet.}
    \item Therapeutic Guidelines (Australia). \textit{Neurological disorders.} https://www.tg.org.au/
    \item Royal Australian College of General Practitioners (RACGP). \textit{Myasthenia gravis resources.}
\end{itemize}

\clearpage
